\documentclass[10pt,a4paper]{amsart}
\usepackage[margin=19mm]{geometry}
\usepackage{amsmath,amssymb,mathtools}
\usepackage[scr=rsfs,cal=euler]{mathalfa}
\usepackage{xfrac}
\usepackage[unicode,hidelinks,bookmarks=false]{hyperref}
\newcommand{\ie}{i.e.\ }
\newcommand{\cf}{cf.\ }
\newcommand{\resp}{respectively\ }
\newcommand{\Subsection}{\subsection}
\providecommand{\nobreakdash}{\nobreak\hbox{-}}
\setlength{\emergencystretch}{2em}
\multlinegap0pt
\usepackage{amssymb}
\newtheorem{theorem}{Theorem}
\newtheorem{lemma}{Lemma}
\newcommand{\PP}{\mathcal{P}}
\newcommand{\T}{\mathcal{T}}
\title{An intermediate conjecture between Goldbach and Dubner:\\ every even number is the sum of a prime and a twin prime}
\author{Tushar Pandey}
\date{Source: arXiv:2608.02381v2 (CC BY 4.0)}
\begin{document}
\maketitle

\noindent\emph{Source and licence.} An intermediate conjecture between Goldbach and Dubner:\\ every even number is the sum of a prime and a twin prime. Version arXiv:2608.02381v2 (CC BY 4.0), distributed by arXiv under the Creative Commons Attribution 4.0 International licence, http://creativecommons.org/licenses/by/4.0/, as recorded in the arXiv licence metadata for this exact version and shown on the abstract page https://arxiv.org/abs/2608.02381v2. Full licence notice: \texttt{licenses/arxiv-2608.02381-CC-BY-4.0.txt}.\par\smallskip

\noindent\textit{Editorial scope.} Counted unit: the article's theorem thm:romanov (source0.tex lines 197--217). The article's complete original proof of it is delivered with it (source0.tex lines 685--766, one \texttt{proof} environment of 82 lines, which the article places away from the statement). No mathematics is written, repaired or re-derived: the statement and the proof are the article's own lines, in the article's own notation, and no retained line is substituted or re-flowed.  Retained uncounted as the source-local context the proof needs: lines 356--359; lines 583--593.  Nothing was omitted from any retained span, so the package carries no omission record.  No retained line contains a citation, so the wrapper carries no bibliography and no bibliography entry.  No retained line contains a figure environment, an image-inclusion command or drawing code, so no image is delivered and the package declares no asset dependency.  Editorial formatting only, in addition: an AMS \texttt{amsart} preamble that restates the article's own declarations and macros used by the retained lines, namely the environments \texttt{theorem}, \texttt{lemma} on separate counters, together with the standard packages those lines need.  The retained environments print in this document as Theorem 1, Lemma 1 in order of appearance, whereas the article prints the same items with the numbering of its own full document.  The complete native source of the article as distributed in the arXiv e-print for this exact version is bundled unchanged as \texttt{sources/206686/source0.tex}.
\begin{equation}\label{eq:twinbound}
  \#\{\, p \le x : p \in \PP,\ p+k \in \PP \,\}
  \;\le\; A\,\mathfrak{S}_2(k)\,\frac{x}{\log^{2} x}.
\end{equation}

\begin{lemma}[The singular series over differences of twin members]
\label{lem:crux}
There are absolute constants $B$ and $z_1$ such that for every $z \ge z_1$,
\[
  \Sigma(z) \;:=\;
  \sum_{\substack{t,\, t' \in \T \cap [1,z] \\ t \ne t'}}
    \mathfrak{S}_2\bigl(|t - t'|\bigr)
  \;\le\; B\,\pi_\T(z)\,\frac{z}{\log^{2} z}.
\]
The statement is unconditional: no hypothesis on $\pi_\T$ is used.
\end{lemma}

\begin{theorem}\label{thm:romanov}
Assume $\mathrm{TW}(c)$, and set
\[
  \mathcal{R} \;=\; \{\, n \text{ even} : n = p + t
    \text{ for some } p \in \PP,\ t \in \T \,\}
  \;=\; (\PP + \T) \cap 2\mathbb{Z}.
\]
Then there is an absolute constant $K$ (independent of $c$, of $z_0$ and of
$\T$) and a threshold $x_0$ depending only on $z_0$, such that
\[
  \#\bigl(\mathcal{R} \cap [1, y]\bigr) \;\ge\; \frac{c}{K}\cdot\frac{y}{2}
  \qquad\text{for all } y \ge 2x_0 .
\]
Consequently $\mathcal{R}$ has lower asymptotic density at least $c/K$ inside
the even integers:
\[
  \liminf_{y \to \infty}
  \frac{\#\{n \le y : n \text{ even},\ n \in \PP + \T\}}{y/2}
  \;\ge\; \frac{c}{K} \;>\; 0 .
\]
\end{theorem}

\begin{proof}[Proof of Theorem~\ref{thm:romanov}]
Let $x$ be a real number, $x \ge \max(z_0, z_1, 3)$, and set
\[
  r_x(n) \;=\; \#\{\, (p,t) \in \PP^{*}\times\T :
    p \le x,\ t \le x,\ p + t = n \,\}.
\]
If $r_x(n) > 0$ then $n$ is a sum of two odd primes, so $n$ is even with
$6 \le n \le 2x$, and $n \in \mathcal{R}$. Put
$\mathcal{N}_x = \{n : r_x(n) > 0\} \subseteq \mathcal{R}\cap[6,2x]$.

\emph{First moment.} Each pair $(p,t) \in (\PP^{*}\cap[1,x]) \times
(\T\cap[1,x])$ is counted once, so
\begin{equation}\label{eq:firstmoment}
  \sum_{n} r_x(n) \;=\; \pi^{*}(x)\,\pi_\T(x).
\end{equation}

\emph{Second moment.} By definition,
\[
  \sum_n r_x(n)^{2} \;=\;
  \#\{\, (p,t,p',t') \in (\PP^{*})^{2}\times\T^{2},\ \text{all} \le x :
    p + t = p' + t' \,\}.
\]
If $p = p'$ then $t = t'$, and conversely; so the quadruples split into a
diagonal part, of size exactly $\pi^{*}(x)\pi_\T(x)$, and an off-diagonal part
with $p \ne p'$ and $t \ne t'$. For an off-diagonal quadruple put
$k = t' - t = p - p' \ne 0$; since $t,t'$ are odd, $k$ is even, and $|k| < x$.
Fix an ordered pair $(t,t')$ with $t \ne t'$. If $k > 0$ the admissible
$(p,p')$ are given by the odd primes $p' \le x$ with $p'+k \le x$ also prime,
and if $k < 0$ by the odd primes $p \le x$ with $p + |k| \le x$ also prime;
either way their number is at most
$\#\{q \le x : q,\, q+|k| \in \PP\}$, which by \eqref{eq:twinbound} is at most
$A\,\mathfrak{S}_2(|k|)\,x/\log^{2}x$. Summing over the ordered pairs $(t,t')$
and applying Lemma~\ref{lem:crux} at $z = x$,
\begin{equation}\label{eq:secondmoment}
  \sum_n r_x(n)^{2}
  \;\le\; \pi^{*}(x)\pi_\T(x)
    + \frac{A\,x}{\log^{2}x}\,\Sigma(x)
  \;\le\; \pi^{*}(x)\pi_\T(x)
    + AB\,\pi_\T(x)\,\frac{x^{2}}{\log^{4}x}.
\end{equation}

\emph{The diagonal is negligible.} By Chebyshev's upper bound,
$\pi(x) \le 2x/\log x$ for $x \ge 2$; indeed the maximum of
$\pi(x)\log x / x$ over $x \ge 2$ is $1.2551\ldots$, attained at $x = 113$.
Hence $\pi^{*}(x) \le 2x/\log x$ and
\[
  \frac{\pi^{*}(x)\pi_\T(x)}{AB\,\pi_\T(x)x^{2}/\log^{4}x}
  \;\le\; \frac{2\log^{3}x}{AB\,x} \;\le\; 1
\]
for $x \ge x_2$, with $x_2$ absolute. For such $x$ the right-hand side
of \eqref{eq:secondmoment} is thus at most
$2AB\,\pi_\T(x)\,x^{2}/\log^{4}x$.

\emph{Cauchy--Schwarz.} From
$\sum_n r_x(n) = \sum_{n \in \mathcal{N}_x} r_x(n)\cdot 1$ we get
$\bigl(\sum_n r_x(n)\bigr)^{2} \le |\mathcal{N}_x| \sum_n r_x(n)^{2}$, so by
\eqref{eq:firstmoment} and the previous paragraph, for
$x \ge \max(z_0, z_1, x_2)$,
\[
  |\mathcal{N}_x| \;\ge\;
  \frac{\pi^{*}(x)^{2}\pi_\T(x)^{2}}{2AB\,\pi_\T(x)\,x^{2}/\log^{4}x}
  \;=\; \frac{\pi^{*}(x)^{2}\,\pi_\T(x)\,\log^{4}x}{2AB\,x^{2}} .
\]
By Chebyshev's lower bound (or $\pi(x) > x/\log x$ for $x \ge 17$) there is an
absolute $x_3$ with $\pi^{*}(x) \ge \tfrac12 x/\log x$ for $x \ge x_3$, whence
$\pi^{*}(x)^{2} \ge x^{2}/(4\log^{2}x)$ and
\[
  |\mathcal{N}_x| \;\ge\; \frac{\pi_\T(x)\log^{2}x}{8AB} .
\]
Finally $\mathrm{TW}(c)$ gives $\pi_\T(x)\log^{2}x \ge c\,x$ for $x \ge z_0$.
So with the absolute constant $K := 8AB = 72 P_1 A^{2}$ and
$x_0 := \max(z_0, z_1, x_2, x_3)$,
\[
  \#\bigl(\mathcal{R}\cap[1,2x]\bigr) \;\ge\; |\mathcal{N}_x|
  \;\ge\; \frac{c}{K}\,x \qquad (x \ge x_0).
\]
Given $y \ge 2x_0$, apply this with $x = y/2$ (all definitions above make
sense for real $x$) to get
$\#(\mathcal{R}\cap[1,y]) \ge \frac{c}{K}\cdot\frac{y}{2}$. Since the number of
even integers in $[1,y]$ is $\lfloor y/2\rfloor$, dividing and letting
$y \to \infty$ gives the density statement.
\end{proof}

\end{document}
